% Regression: a wire declares its stroke.  A rail is solid unless it says
% otherwise; a dashed rail carries an index the picture draws but does not
% contract, and a dotted one carries a lattice lying under the sheet being
% drawn.  The pattern is a named metric, so the ink names a house style and
% never a dash length of its own, and the weight and colour of the rail are
% untouched.
\documentclass{standalone}
\usepackage{tenkz}
\makeatletter
\ExplSyntaxOn
\file_input:n { tenkz-kernel.code.tex }
\int_new:N \g__tenkz_test_stroke_solid_int
\int_new:N \g__tenkz_test_stroke_dashed_int
\int_new:N \g__tenkz_test_stroke_dotted_int
\cs_new_eq:NN
  \__tenkz_test_wire_stroke:nn
  \__tenkz_kernel_r_wire_stroke:nn
\cs_set_protected:Npn \__tenkz_kernel_r_wire_stroke:nn #1#2
  {
    \exp_args:NV \__tenkz_model_get:nnN
      \l__tenkz_kernel_r_wire_record_tl {stroke} \l_tmpa_tl
    \quark_if_no_value:NTF \l_tmpa_tl
      {
        \tl_if_in:NnT \l__tenkz_kernel_r_wire_ink_tl {tenkz~wire~}
          {\errmessage{an~undeclared~rail~acquired~a~dash~pattern}}
        \int_gincr:N \g__tenkz_test_stroke_solid_int
      }
      {
        \str_case:Vn \l_tmpa_tl
          {
            {solid}
              {
                \tl_if_in:NnT \l__tenkz_kernel_r_wire_ink_tl {tenkz~wire~}
                  {\errmessage{a~solid~rail~acquired~a~dash~pattern}}
                \int_gincr:N \g__tenkz_test_stroke_solid_int
              }
            {dashed}
              {
                \tl_if_in:NnF \l__tenkz_kernel_r_wire_ink_tl
                  {tenkz~wire~dashed}
                  {\errmessage{a~dashed~rail~lost~its~dash~pattern}}
                \int_gincr:N \g__tenkz_test_stroke_dashed_int
              }
            {dotted}
              {
                \tl_if_in:NnF \l__tenkz_kernel_r_wire_ink_tl
                  {tenkz~wire~dotted}
                  {\errmessage{a~dotted~rail~lost~its~dot~pattern}}
                \int_gincr:N \g__tenkz_test_stroke_dotted_int
              }
          }
      }
    \__tenkz_test_wire_stroke:nn {#1} {#2}
  }
\AtEndDocument
  {
    \int_compare:nNnF { \g__tenkz_test_stroke_solid_int } = {2}
      {\errmessage{solid~rail~count~changed}}
    \int_compare:nNnF { \g__tenkz_test_stroke_dashed_int } = {2}
      {\errmessage{dashed~rail~count~changed}}
    \int_compare:nNnF { \g__tenkz_test_stroke_dotted_int } = {2}
      {\errmessage{dotted~rail~count~changed}}
  }
\ExplSyntaxOff
\makeatother
\begin{document}
\tenkzkernel
% Three contracted rails: the default, the dashed spelling, the dotted one.
\begin{tenkz}[rows={wire,wire,wire}, cols=3, bonds=none]
  \tn[at=(1,1), name=a, skin=dot, ports={0:virtual}]{}
  \tn[at=(1,3), name=b, skin=dot, ports={180:virtual}]{}
  \tn[at=(2,1), name=c, skin=dot, ports={0:virtual}]{}
  \tn[at=(2,3), name=d, skin=dot, ports={180:virtual}]{}
  \tn[at=(3,1), name=e, skin=dot, ports={0:virtual}]{}
  \tn[at=(3,3), name=f, skin=dot, ports={180:virtual}]{}
  \tnwire[name=plain]{a.0}{b.180}
  \tnwire[stroke=dashed, name=dash]{c.0}{d.180}
  \tnwire[stroke=dotted, name=dot]{e.0}{f.180}
\end{tenkz}
% The spelling rides an open leg and a diagonal as readily as a straight
% contracted rail, and solid= is the default said out loud.
\begin{tenkz}[rows={wire,wire,wire}, cols=3, bonds=none]
  \tn[at=(3,1), name=g, skin=dot, ports={45:virtual, 270:virtual}]{}
  \tn[at=(1,3), name=h, skin=dot, ports={225:virtual}]{}
  \tn[at=(2,2), name=k, skin=dot, ports={180:virtual}]{}
  \tnwire[stroke=dotted, name=diagonal]{g.45}{h.225}
  \tnwire[stroke=dashed, name=pendant]{k.180}{open w}
  \tnwire[stroke=solid, name=said]{g.270}{open s}
\end{tenkz}
\end{document}
